\section{PRODUCTION-HARDENING LESSONS}\label{sec:lessons}

The central finding of Section~\ref{sec:intro} has a second, less statistical face. This system had
several defects invisible to its own test suite for the same reason the mock evaluation flatters its
headline numbers: the path exercised was not the path production takes. Each lesson below surfaced
through an audit or an experiment rather than a failing test, was measured before it was fixed, and
is documented in the accompanying artifact.

\subsection{L1: The Shared Resource Nobody Declared Concurrency-Unsafe}
The base design asserts that execution is deterministic and auditable. Versioned partition creation,
though, read the current maximum version and then dropped, created, and registered a table across
separate statements with no lock between them, so two concurrent callers targeting the same
partition could compute the same next version and race. The defect showed up as an intermittent
continuous-integration failure and was root-caused rather than retried away: a recovery-agent replay
had triggered the same DAG as an unrelated test. The fix serializes the read-modify-write behind a
database advisory lock. \emph{Lesson:} once an agent can trigger real work, every shared resource it
touches becomes a concurrency surface that a single-actor test never exercises.

\subsection{L2: An Action Could Be Lost From the Audit Trail}
The audit row was originally written \emph{after} execution, so a crash, an out-of-memory kill, or a
database blip in that window could leave an action that genuinely happened with no record of it, and
because the control loop swallows per-tick exceptions, the failure stayed silent. For a system whose
central claim is that every action is policy-gated and auditable, an executed action with no audit
trail falsifies that claim outright. The intent row, carrying the policy verdict already known, is
now committed \emph{before} execution and updated afterward with the outcome. \emph{Lesson:} record
the intent before the action, not after.

\subsection{L3: The Detector That Had No Callers}
The table that both recovery-time statistics and decision scoring depend on was written only by the
fault injector. The monitoring agent updated an existing row's detection time but never created a new
one, leaving a complete, unit-tested, deterministic anomaly detector with no caller at all. In a
deployment without the injector, a genuinely detected incident would have reached the audit table but
never become an incident, so production could not have reproduced the experiment's own metrics. The
gap turned up while planning a later feature, by reading the code rather than through a failing
test. \emph{Lesson:} a component can pass every test it owns and still be disconnected, so trace the
data path from production input all the way to the reported metric.

\subsection{L4: Two Taxonomies That Never Met}
Decision correctness is scored against an accepted-mitigation set (Eq.~\ref{eq:decision}) keyed by the \emph{injected
scenario}'s name, while live incidents carry different fault-type names altogether. Applying that set
to live incidents would have marked every real incident ``incorrect by construction.'' A mapping from
live fault types onto the same accepted-mitigation set was added before the metric was ever exposed
outside the experiment harness. \emph{Lesson:} a metric defined in the experiment's own vocabulary
cannot move into production without an explicit mapping.

\subsection{L5: Hermeticity by Enumeration Does Not Scale}
One class of unit-test failure kept recurring: a test that forgot to mock one module's database
handle turned a millisecond test into a thirty-second pool timeout. We patched it instance by
instance four separate times before addressing the underlying mechanism. Every such call now runs
through a single connection-pool accessor patched to fail fast under test. Mutation-testing the guard
by removing one load-bearing mock failed six tests within seconds, each naming the exact cause.
\emph{Lesson:} once the same bug surfaces a fourth time, fix the mechanism, not the instance.

\subsection{L6: What the Adversarial Corpus Found}
Grading a generated corpus against an independent oracle surfaced three things (numbers in
Section~\ref{sec:results}). First, the cost policy priced a scaling target against the budget but
said nothing about its floor: zero workers, which halts ingestion entirely, was budget-legal, and so
was a negative ``scale-down.'' Non-numeric targets raised inside the gate before the audit row was
written, so a rejected proposal could leave no record at all. The fix lives in the contract layer,
where scaling targets must now be integers of at least one, with any violation taking the ordinary
invalid-output path.

Second, under a newer policy-engine release than the one this project pins, some float budget
comparisons evaluated as within budget when they were not, while the pinned release stayed correct.
We report the observation without root-causing it or filing it upstream, since the consequence holds
either way: the pin is load-bearing, and the corpus now serves as an integration test that would
catch a version bump that changes these semantics again.

Third, once the contract fix landed, the policy agreed with the independent specification on every
graded case. The containment claim is therefore a measured property with a confidence interval, not
an assertion repeated from the code under test. \emph{Lesson:} a safety property a policy's own
authors assert is not evidence until an independent oracle has tried to break it.

\subsection{L7: The Pilot Found a Dead Model}
Before the live campaign, an eight-run pilot at the real paper timings measured cost and wall time.
Within minutes it found that the provider had retired the small model configured for the monitoring
agent (HTTP 410, end of life). Since monitoring is the detector in every non-baseline configuration,
the full campaign would have run against a detector degraded to a no-op, with every downstream number
measuring provider retirement rather than the system itself. The supervisor's own degraded-streak
guard would eventually have caught this too, but only after several wasted runs, while the pilot
caught it on the very first call. \emph{Lesson:} the cheapest validity check available in a long,
paid experiment is a short run at the real settings, since a model listed by a provider is not
necessarily a model that still answers.
